\documentclass{article}
\usepackage{amsmath, amssymb, amsthm}
\usepackage{hyperref}
\usepackage{geometry}
\geometry{margin=1in}

\title{Erd\H{o}s Problem \#951}
\date{Last updated: 2025-08-31}
\author{Source: erdosproblems.com}

\begin{document}
\maketitle

\noindent\textbf{Status:} OPEN \quad \textbf{No prize}

\noindent\textbf{Tags:} number theory

\medskip
\noindent\textbf{Problem Statement:}

Let $1<a_1<\cdots$ be a sequence of real numbers such that\[\left\lvert \prod_i a_i^{k_i}-\prod_j a_j^{\ell_j}\right\rvert \geq 1\]for every distinct pair of non-negative finitely supported integer tuples $k_i,\ell_j\geq 0$. Is it true that\[\#\{ a_i \leq x\} \leq \pi(x)?\]

\bigskip
\noindent\textbf{Remarks:}

Erd\H{o}s \cite{Er77c} said this question was asked 'during [his] lecture at Queens College [by] one member of the audience (perhaps S. Shapiro)'. (Although in \cite{Er80} he seems sure it was Shapiro, but also recalled that he had himself asked this in \cite{Er69}.) In \cite{Er80} he further asks whether equality holds if and only if the $a_i$ are the set of primes.

Such a sequence of $a_i$ is sometimes called a set of Beurling prime numbers (and the sequence of products called the associated generalised integers).

Beurling conjectured that if the number of reals in $[1,x]$ of the form $\prod a_i^{k_i}$ is $x+o(\log x)$ then the $a_i$ must be the sequence of primes.

It is unclear whether this is intended to hold for all $x$ or just all sufficiently large $x$. The former question can be disproved by finite calculation (since any finite sequence of $a_i$ can be extended to a suitable infinite sequence via a greedy algorithm). A finite counterexample for $x=10$ was found by ChatGPT-5.2 Pro (prompted by Leeham).

\bigskip
\noindent\textbf{References:}

[Er69] Erd\H{o}s, Paul, Some applications of graph theory to number theory. The Many Facets of Graph Theory (Proc. Conf.,
Western Mich. Univ., Kalamazoo, Mich., 1968) (1969), 77-82.
 
 [Er77c] Erd\H{o}s, Paul, Problems and results on combinatorial number theory. III. Number theory day (Proc. Conf., Rockefeller Univ.,
New York, 1976) (1977), 43-72.
 
 [Er80] Erd\H{o}s, Paul, A survey of problems in combinatorial number theory. Ann. Discrete Math. (1980), 89-115.

\bigskip
\noindent\small{Source: \url{https://www.erdosproblems.com/951}}
\end{document}
